\subsection{无穷小自同构}

引理 1. 设 G 是李群，\(\alpha\) 是 G 上的向量场。对所有 \(g \in \mathbf{G}\)，令

\[
\beta (g) = \alpha (g) g ^ {- 1} \in \mathrm{L} (\mathrm{G}).
\]

以下条件等价：

\begin{enumerate}
\def\labelenumi{(\roman{enumi})}
\item
  \(\alpha\) 是从群 G 到群 T(G) 的同态；
\item
  对所有 \(g, g'\) 属于 \(\mathbf{G}\)，都有 \(\alpha(gg') = \alpha(g)g' + g\alpha(g')\)；
\item
  对所有 \(g, g'\) 属于 \(\mathbf{G}\)，都有 \(\beta(gg') = \beta(g) + (\mathrm{Ad}g)\beta(g')\)。
\end{enumerate}

条件 (i) 意味着，对所有 \(g, g'\) 属于 \(G\)，在群 \(T(G)\) 中有：

\[
\beta (g) g \beta (g ^ {\prime}) g ^ {\prime} = \beta (g g ^ {\prime}) g g ^ {\prime}
\]

或者

\[
\beta (g) ((\mathrm{Ad} g) \beta (g ^ {\prime})) g g ^ {\prime} = \beta (g g ^ {\prime}) g g ^ {\prime}.
\]

但是 \(\beta(g)\) 与 (Ad \(g)\beta(g')\) 在 T(G) 中的乘积，正是 \(\beta(g)\) 与 (Ad \(g)\beta(g')\) 在 L(G) 中的和（第 2 节第 1 号命题 2）。因此 (i) \(\Leftrightarrow\) (iii)。另一方面，条件 (ii) 可写成 \(\beta(gg')gg' = \beta(g)gg' + g\beta(g')g'\)，或者

\[
\beta (g g ^ {\prime}) = \beta (g) + (\mathrm{Ad} g) \beta (g ^ {\prime})
\]

从而 (ii) \(\Leftrightarrow\) (iii)。

定义 1. 设 G 是李群。G 的无穷小自同构是 G 上满足引理 1 条件的任一解析向量场。

引理 2. 设 \(\mathbf{K}'\) 是 \(\mathbf{K}\) 的非离散闭子域，A 是 \(\mathbf{K}'\)-流形，B 和 C 是 K-流形，且 \(f\) 是 \(\mathbf{K}'\)-解析映射，从 \(\mathbf{A} \times \mathbf{B}\) 到 C。假设对所有 \(a \in \mathbf{A}\)，映射 \(b \mapsto f(a, b)\) 从 \(\mathbf{B}\) 到 \(\mathbf{C}\) 都是 K-解析的。于是对所有 \(t \in \mathbf{TA}\)，从 TB 到 TC 的映射 \(u \mapsto (\mathrm{Tf})(t, u)\) 是 K-解析的。

固定 \(t \in \mathrm{TA}\)，记 \(g(u) = (\mathrm{Tf})(t,u)\)。显然，\(g\) 是 K'-解析的。由《可微与解析流形》R，5.14.6，只需证明 \(g\) 的切映射是 K-线性的。可以假定 A、B、C 分别是 K'、K、K 上完备可赋范空间 E、F、G 中 0 的开邻域，且 \(t\) 是 A 在 0 处的切向量。将 TA、TB、TC 分别与 A × E、B × F、C × G 识别，并将 \(t\) 与 E 的一个元素识别。于是对所有 \((x,y) \in \mathrm{TB} = \mathrm{B} \times \mathrm{F}\)。

\[
g (x, y) = \left(f (0, x), \left(\mathrm{D} _ {1} f\right) (0, x) (t) + \left(\mathrm{D} _ {2} f\right) (0, x) (y)\right).
\]

将 \(\mathrm{T(B\times F)}\) 与 \((\mathbf{B}\times \mathbf{F})\times (\mathbf{F}\times \mathbf{F})\) 识别，将 \(\mathrm{T(C\times G)}\) 与 \((\mathbf{C}\times \mathbf{G})\times (\mathbf{G}\times \mathbf{G})\) 识别。于是，对所有

\[
((x, y), (h, k)) \in \mathrm{T} (\mathrm{B} \times \mathrm{F}) = (\mathrm{B} \times \mathrm{F}) \times (\mathrm{F} \times \mathrm{F}),
\]

\((\mathrm{Tg})((x,y),(h,k)) = ((a,b),(c,d)),\) 其中

\[
\begin{array}{r l} & {a = f (0, x),} \\ & {b = (\mathrm{D} _ {1} f) (0, x) (t) + (\mathrm{D} _ {2} f) (0, x) (y), \qquad c = (\mathrm{D} _ {2} f) (0, x) (h),} \\ & {d = (\mathrm{D} _ {2} \mathrm{D} _ {1} f) (0, x) (t, h) + (\mathrm{D} _ {2} \mathrm{D} _ {2} f) (0, x) (y, h) + (\mathrm{D} _ {2} f) (0, x) (k).} \end{array}
\]

现在固定 \((x,y)\in \mathbf{B}\times \mathbf{F}\)。我们需要证明映射 \((h,k)\mapsto (c,d)\) 从 \(\mathbf{F}\times \mathbf{F}\) 到 \(\mathbf{G}\times \mathbf{G}\) 是 K-线性的。由于从 B 到 C 的映射 \(x\mapsto f(0,x)\) 是 K-解析的，映射

\[
\begin{array}{c} (h, k) \mapsto (\mathrm{D} _ {2} f) (0, x) (h), \qquad (h, k) \mapsto (\mathrm{D} _ {2} \mathrm{D} _ {2} f) (0, x) (y, h), \\ (h, k) \mapsto (\mathrm{D} _ {2} f) (0, x) (k) \end{array}
\]

是 K-线性的。另一方面

\[
\left(\mathrm{D} _ {2} \mathrm{D} _ {1} f\right) (0, x) (t, h) = \lim _ {\lambda \in K ^ {*}, \lambda \rightarrow 0} \lambda^ {- 1} \left(\left(\mathrm{D} _ {2} f\right) (\lambda t, x) (h) - \left(\mathrm{D} _ {2} f\right) (0, x) (h)\right)
\]

并且，对固定的 \(\lambda\)，映射 \(x \mapsto f(\lambda t, x)\) 是 K-解析的，所以映射 \(h \mapsto (\mathrm{D}_2f)(\lambda t, x)(h)\) 是 K-线性的。

命题 1. 设 \(\mathbf{K}'\) 是 \(\mathbf{K}\) 的非离散闭子域，\(\mathbf{G}\) 是 \(\mathbf{K}\) 上的李群，\(\mathbf{V}\) 是 \(\mathbf{K}'\) 上的流形，且 \((v, g) \mapsto vg\) 是 \(\mathbf{K}'\)-解析映射，从 \(\mathbf{V} \times \mathbf{G}\) 到 \(\mathbf{G}\)。假设对所有 \(v \in \mathbf{V}\)，映射 \(g \mapsto vg\) 从 \(\mathbf{G}\) 到 \(\mathbf{G}\) 都是 \(\mathbf{G}\) 的自同构。设 \(\varepsilon\) 是 \(\mathbf{V}\) 的一个元素，满足 \(\varepsilon g = g\) 对所有 \(g \in \mathbf{G}\) 成立，且 \(a \in T_{\varepsilon}(\mathbf{V})\)。则 \(g \mapsto ag\) 是 \(\mathbf{G}\) 上的向量场，也是 \(\mathbf{G}\) 的无穷小自同构。

对 \(v \in \mathrm{V}\)、\(g_1 \in \mathrm{G}\)、\(g_2 \in \mathrm{G}\)，有 \(v(g_1g_2) = (vg_1)(vg_2)\)。因此，对 \(u_1 \in \mathrm{TG}\)、\(u_2 \in \mathrm{TG}\)，有 \(a(u_1u_2) = (au_1)(au_2)\)（第 2 节第 1 号命题 3）。特别地，映射 \(g \mapsto ag\) 从 \(G\) 到 \(TG\) 是群同态。另一方面，由引理 2，这个映射是解析的。

命题 2. 设 G 是实或复李群，\(\alpha\) 是 G 的无穷小自同构。存在 K 在 G 上的解析运算律 \((\lambda, g) \mapsto \phi_{\lambda}(g)\)，具有以下性质：

\begin{enumerate}
\def\labelenumi{(\arabic{enumi})}
\item
  若 D 是相应的无穷小运算律，则 D(1) = α；
\item
  对所有 \(\lambda \in \mathbf{K}\)，\(\phi_{\lambda} \in \operatorname{Aut} G\)。
\end{enumerate}

\begin{enumerate}
\def\labelenumi{(\alph{enumi})}
\item
  对所有 \(\mu > 0\)，令 \(\mathbf{K}_{\mu}\) 为半径为 \(\mu\)、位于 \(\mathbf{K}\) 中的开球。对所有 \(g \in G\)，令 \(\mathcal{F}_g\) 为由 \(f\)（\(\alpha\) 的、在 \(\mathbf{K}_{\mu}\) 上定义的解析积分曲线）组成的集合，且满足 \(f(0) = g\)。由《可微与解析流形》R，9.1.3 和 9.1.5，\(\mathcal{F}_g\) 非空，且 \(\mathcal{F}_g\) 的两个元素在其定义域的交集上重合；令 \(\mu(g)\) 为满足以下条件的 \(\mu\)：\(\mathcal{F}_g\) 中存在定义在 \(\mathbf{K}_{\mu}\) 上的元素；\(\mathcal{F}_g\) 中存在唯一一个定义在 \(\mathbf{K}_{\mu(0)}\) 上的元素；记其为 \(f_g\)。
\item
  设 \(g_1, g_2\) 属于 \(\mathbf{G}\)，\(f_1 \in \mathcal{F}_{g_1}\)，\(f_2 \in \mathcal{F}_{g_2}\)，且 \(f_1\) 与 \(f_2\) 定义在同一个球 \(\mathbf{K}_\mu\) 上。则 \(f_1f_2: \mathbf{K}_\mu \to \mathbf{G}\) 是解析的，且 \((f_1f_2)(0) = g_1g_2\)。另一方面，对所有 \(\lambda \in \mathbf{K}_\mu\)，
\end{enumerate}

\[
\begin{array}{r l} (\mathrm{T} _ {\lambda} (f _ {1} f _ {2})) 1 & = (\mathrm{T} _ {\lambda} f _ {1}) 1 \cdot f _ {2} (\lambda) + f _ {1} (\lambda) \cdot (\mathrm{T} _ {\lambda} f _ {2}) 1 \quad (\S 2, \text { Proposition   7 }) \\ & = \alpha (f _ {1} (\lambda)) f _ {2} (\lambda) + f _ {1} (\lambda) \alpha (f _ {2} (\lambda)) \\ & = \alpha ((f _ {1} f _ {2}) (\lambda)) \quad \text {(Lemma   1)} \end{array}
\]

从而 \(f_{1}f_{2}\in \mathcal{F}_{g_{1}g_{2}}\)。这就证明了 \(\mu (g_1g_2)\geqslant \inf (\mu (g_1),\mu (g_2))\)。

\begin{enumerate}
\def\labelenumi{(\alph{enumi})}
\setcounter{enumi}{2}
\tightlist
\item
  由《可微与解析流形》R，9.1.4 和 9.1.5，存在 G 中 \(e\) 的一个邻域 V，使得 \(\sigma = \inf_{g \in \mathbf{V}} \mu(g) > 0\)。令 \(h \in \mathbf{G}\)，C 为其连通分支。由 (b)，对所有 \(h' \in \mathbf{C}\)，\(\mu(h') \geqslant \inf(\sigma, \mu(h)) > 0\)。另一方面，函数 \(f_h'\)（其中 \(h' \in \mathbf{C}\)）的值都在 C 中。由《可微与解析流形》R，9.1.4 和 9.1.5，C 中 \(\mu = +\infty\)，最终 G 中 \(\mu = +\infty\)。现在令 \(f_g(\lambda) = \phi_{\lambda}(g)\)，其中 \(g \in \mathbf{G}\) 和 \(\lambda \in \mathbf{K}\) 任意。由《可微与解析流形》R，9.1.4 和 9.1.5，映射 \((\lambda, g) \mapsto \phi_{\lambda}(g)\) 是 K 在 G 上的解析运算律。显然，若 D 是相应的无穷小运算律，则 D(1) = \(\alpha\)。由 (b)，
\end{enumerate}

\[
\phi_ {\lambda} (g _ {1} g _ {2}) = \phi_ {\lambda} (g _ {1}) \phi_ {\lambda} (g _ {2})
\]

对所有 \(\lambda \in \mathbf{K}\)、\(g_{1} \in \mathbf{G}\)、\(g_{2} \in \mathbf{G}\)，均成立。

命题 3. 假设 K 是超度量的。设 G 是紧李群，\(\alpha\) 是 G 的无穷小自同构。存在 K 的一个开子群 I 以及 I 在 G 上的解析运算律 \((\lambda, g) \mapsto \phi_{\lambda}(g)\)，具有以下性质：

\begin{enumerate}
\def\labelenumi{(\arabic{enumi})}
\item
  若 D 是相应的无穷小运算律，则 D(1) = α；
\item
  对所有 \(\lambda \in \mathbf{I}\)，\(\phi_{\lambda} \in \operatorname{Aut} G\)。
\end{enumerate}

由于 G 是紧的，存在 K 的一个开子群 I' 以及 I' 在 G 上的解析运算律 \((\lambda, g) \mapsto \phi_{\lambda}(g)\)，具有本命题的性质 (1)（第 4 节第 7 号定理 6 的推论 2）。记 \(\phi_{\lambda}(g) = f_{g}(\lambda)\)，其中 \(\lambda \in I'\) 和 \(g \in G\)。于是，对 \(g_{1}, g_{2}\) 属于 \(G\) 和 \(\lambda \in I'\)，

\[
\begin{array}{r l} (\mathrm{T} _ {\lambda} (f _ {g _ {1}} f _ {g _ {2}})) 1 & = (\mathrm{T} _ {\lambda} f _ {g _ {1}}) 1 \cdot f _ {g _ {2}} (\lambda) + f _ {g _ {1}} (\lambda) \cdot (\mathrm{T} _ {\lambda} f _ {g _ {2}}) 1 \\ & = \alpha (f _ {g _ {1}} (\lambda)) f _ {g _ {2}} (\lambda) + f _ {g _ {1}} (\lambda) \alpha (f _ {g _ {2}} (\lambda)) \\ & = \alpha (f _ {g _ {1}} (\lambda) f _ {g _ {2}} (\lambda)) \end{array}
\]

并且 \((f_{g_1}f_{g_2})(0) = g_1g_2 = f_{g_1g_2}(0)\)。因此 \(f_{g_1'g_2'}(\lambda) = f_{g_1'}(\lambda)f_{g_2'}(\lambda)\) 对

\[
(g _ {1} ^ {\prime}, g _ {2} ^ {\prime}, \lambda)
\]

在 \((g_1, g_2, 0)\) 的一个邻域内成立（《可微与解析流形》R，9.1.8）。由于 G 是紧的，存在 I' 的一个开子群 I，使得 \(f_{g_1 g_2}(\lambda) = f_{g_1}(\lambda) f_{g_2}(\lambda)\) 对所有 \(g_1 \in G\)、\(g_2 \in G\)、\(\lambda \in I\) 成立。换言之，\(\phi_\lambda \in \operatorname{Aut} G\)（当 \(\lambda \in I\) 时）。

引理 3. 设 G 和 G' 是李群，\(\phi\) 是从 G 到 \(\operatorname{Aut}(\mathbf{G}')\) 的同态。令 \(f(g, g') = (\phi(g))(g')\)，其中 \(g \in \mathbf{G}\)、\(g' \in \mathbf{G}'\)。考虑以下条件：

\begin{enumerate}
\def\labelenumi{(\roman{enumi})}
\item
  \(f\) 是解析的；
\item
  \(f\) 在 \((e_{\mathbf{G}}, e_{\mathbf{G}^{\prime}})\) 的一个邻域内是解析的；
\item
  对所有 \(g' \in G'\)，映射 \(g \mapsto f(g, g')\) 都是解析的。
\end{enumerate}

则 (i) \(\Leftrightarrow\) ((ii) 且 (iii))。若 G 是连通的，则 (i) \(\Leftrightarrow\) (ii)。

显然 (i) 蕴含 (ii) 和 (iii)。设 \(g_0 \in G, g_0' \in G'\)。对所有 \(g \in G, g' \in G'\)，

\[
f \left(g g _ {0}, g ^ {\prime} g _ {0} ^ {\prime}\right) = (\phi (g) \phi (g _ {0})) \left(g ^ {\prime} g _ {0} ^ {\prime}\right) = \phi (g) \left(\phi (g _ {0}) g ^ {\prime}\right)\cdot \phi (g) \left(\phi (g _ {0}) g _ {0} ^ {\prime}\right).
\]

这证明了蕴含 ((ii) 且 (iii)) \(\Rightarrow\) (i)。最后，若 \(\mathbf{G}'\) 是连通的，则 \(\mathbf{G}'\) 由 \(e_{\mathbf{G}'}\) 的每个邻域生成，从而 (ii) \(\Rightarrow\) (iii)。

